\section{Rapid Eye Movement (REM) Sleep}\label{sec:rem}

First discovered in 1953, REM sleep is a neurobiological paradox. While the body is almost completely paralyzed (a state known as REM atonia, which prevents us from acting out our dreams), the brain becomes highly active, with EEG patterns resembling an awake, alert state~\cite{fuller2006}. This is the stage most associated with vivid, bizarre, and narrative-driven dreaming.

Key characteristics of REM sleep include:
\begin{itemize}
    \item \textbf{Rapid Eye Movements:} The darting eye movements that give this stage its name, often interpreted as correlates of visual dream imagery~\cite{fuller2006}.
    \item \textbf{Muscle Atonia:} A protective paralysis of voluntary muscles.
    \item \textbf{Activated Brain:} Increased cerebral blood flow and oxygen consumption, particularly in emotional centers like the amygdala and memory centers like the hippocampus~\cite{fuller2006}.
    \item \textbf{Physiological Fluctuation:} Irregular breathing, heart rate, and body temperature.
\end{itemize}
